\section{Experiments and Results}
\label{sec:exp}

\paragraph{Dataset.}
We use the Amazon Reviews 2023 Books subset~\cite{hou2024amazon}
(Table~\ref{tab:dataset}).
The raw corpus has 29.5M reviews and 4.4M items.
We retain users with at least 5 interactions: 100k users over a 3.08M-item
catalog. 375k items (12.2\%) carry a co-purchase edge and are covered by the
behavioral graph embedding; average sequence length is 33.5 clicks.

\input{tables/dataset_stats}

Three dataset properties inform system design.
Item cover images exist for 99.1\% of items but user review images for only
1.8\%, motivating item-side visual embeddings.
About 80\% of reviews are 4--5 stars~\cite{hou2024amazon}, so ratings cannot
serve as a primary relevance signal.
Review text length is right-skewed (median 37, mean 86 words), fitting
standard transformer budgets without aggressive truncation.

The last interaction of each user is the test item; remaining history forms
the training sequence.

\paragraph{Evaluation Metrics.}
We report HR@10 and NDCG@10.
Following the protocol of SASRec~\cite{kang2018self} and
DIF-SASRec~\cite{xie2022difsr}, each held-out item is ranked against 99 random
negatives (random-chance HR@10 = 0.10).
Negatives are sampled uniformly at random from items absent from each user's
training history.
We also evaluate under a 999-negative protocol (random-chance HR@10 = 0.01)
to assess robustness (Section~\ref{ssec:robustness}).

\paragraph{Baselines.}
% Called via % Called via % Called via \input{tables/baselines} inside the Baselines paragraph.
We compare against the following baselines and ablation variants:
\begin{description}
  \item[\textbf{Content-KNN}] cosine nearest-neighbor retrieval on BGE-M3
        embeddings, without sequential modeling.
  \item[\textbf{DIF-SASRec~\cite{xie2022difsr} (Pipeline~B)}] DIF-SASRec with
        BGE-M3 content embeddings, category auxiliary loss, and online
        adaptation (Pipeline~B of our system).
  \item[\textbf{Pipeline~A (Cleora~\cite{cleora2021})}] Cleora behavioral graph
        embeddings with BGE-M3 profile similarity ranking, without DIF-SASRec
        scoring.
  \item[\textbf{System (A $\cup$ B)}] the full dual-pipeline system, returning
        the union of the top-10 lists from Pipeline~A and Pipeline~B.
\end{description}
 inside the Baselines paragraph.
We compare against the following baselines and ablation variants:
\begin{description}
  \item[\textbf{Content-KNN}] cosine nearest-neighbor retrieval on BGE-M3
        embeddings, without sequential modeling.
  \item[\textbf{DIF-SASRec~\cite{xie2022difsr} (Pipeline~B)}] DIF-SASRec with
        BGE-M3 content embeddings, category auxiliary loss, and online
        adaptation (Pipeline~B of our system).
  \item[\textbf{Pipeline~A (Cleora~\cite{cleora2021})}] Cleora behavioral graph
        embeddings with BGE-M3 profile similarity ranking, without DIF-SASRec
        scoring.
  \item[\textbf{System (A $\cup$ B)}] the full dual-pipeline system, returning
        the union of the top-10 lists from Pipeline~A and Pipeline~B.
\end{description}
 inside the Baselines paragraph.
We compare against the following baselines and ablation variants:
\begin{description}
  \item[\textbf{Content-KNN}] cosine nearest-neighbor retrieval on BGE-M3
        embeddings, without sequential modeling.
  \item[\textbf{DIF-SASRec~\cite{xie2022difsr} (Pipeline~B)}] DIF-SASRec with
        BGE-M3 content embeddings, category auxiliary loss, and online
        adaptation (Pipeline~B of our system).
  \item[\textbf{Pipeline~A (Cleora~\cite{cleora2021})}] Cleora behavioral graph
        embeddings with BGE-M3 profile similarity ranking, without DIF-SASRec
        scoring.
  \item[\textbf{System (A $\cup$ B)}] the full dual-pipeline system, returning
        the union of the top-10 lists from Pipeline~A and Pipeline~B.
\end{description}


\paragraph{Implementation Details.}
All experiments run on a single RTX~5060~Ti (16\,GB VRAM).
DIF-SASRec is trained with AdamW, sampled softmax (512 negatives),
and an auxiliary category loss.
FAISS HNSW~\cite{malkov2018efficient} handles nearest-neighbor search;
Cleora uses 3 Markov propagation steps.
BGE-M3 and CLIP inference run in FP16.
